\section{总结与延伸}

Week 4 从两个互补视角探讨了 Coding Agent 的使用模式：Mihail Eric 讲解了 Agent Manager 的方法论和管理技术，Boris Cherny 则从 Claude Code 创造者的角度分享了工具设计哲学和实践经验。

\subsection{关键要点}

\begin{itemize}
    \item 软件开发正从``人管人''转向``人管 Agent''——这是第四次组织形态转型
    \item 四种 Agent 指导技术：行为文件、hooks、commands、subagents
    \item 关键实践：防护栏（tests + CI/CD）、审计标记、模型匹配、频繁 checkpoint
    \item Claude Code 设计哲学：terminal-native、low-level access、infinitely hackable
    \item 根据任务类型选择工作流：explore-plan-code、TDD、视觉迭代
    \item 为 6 个月后的模型能力做准备
\end{itemize}

\subsection{落地清单}

\begin{enumerate}
    \item 为常见任务建立可复用的 commands、hooks 和行为文件
    \item 用 subagents 处理明确分工的并行任务，而不是让所有工作挤进一个上下文
    \item 把测试、CI 和人工审查当作 Agent 工作流的硬边界，而非事后补救
\end{enumerate}

\subsection{角色分工对照表}

\begin{table}[H]
\centering
\begin{tabular}{p{0.22\textwidth} p{0.28\textwidth} p{0.28\textwidth}}
\toprule
角色 & 主要职责 & 失败时的补救方式 \\
\midrule
人类 Agent Manager & 定义目标、拆分任务、设置边界、做最终取舍 & 重新划分任务、补充约束、决定回滚或继续 \\
执行型 Agent & 阅读上下文、修改代码、运行测试、生成提交 & 通过 hooks / tests / review 暴露问题并重试 \\
审查与 CI 系统 & 提供客观验证、阻止明显错误进入主干 & 用标准化检查和可回放日志减少隐性风险 \\
\bottomrule
\end{tabular}
\caption{Agent Manager 模式下的人机职责边界}
\end{table}

\subsection{团队采用成熟度模型}

从课程内容可以提炼出一个很实用的 adoption 路径：

\begin{enumerate}
    \item \textbf{个人试验期}：个别工程师用 Agent 处理问答、脚手架代码和局部重构
    \item \textbf{团队规范期}：开始沉淀 commands、hooks、review 规则和统一的行为文件
    \item \textbf{并行协作期}：Agent 被用于明确分工的多线程任务，开发者开始承担 Agent Manager 角色
    \item \textbf{系统整合期}：Agent 工作流与 CI、代码审查、知识库和项目管理工具深度打通
\end{enumerate}

\begin{knowledgebox}{成熟度提升的标志不是 ``用了更多 Agent''，而是 ``失败更可控''}
真正进入成熟阶段的团队，不是看上去更炫，而是当 Agent 出错时更容易被测试、hooks、代码审查和责任边界及时捕获。这也是为什么 Boris 和 Mihail 都反复强调 guardrails，而不是只强调速度。
\end{knowledgebox}

\subsection{拓展阅读}

\begin{itemize}
    \item Claude Code 官网：\url{https://claude.ai/code}
    \item Claude Code SDK 文档：\url{https://docs.anthropic.com/en/docs/claude-code}
    \item Awesome Claude Agents：\url{https://github.com/vijaythecoder/awesome-claude-agents}
    \item SuperClaude Framework：\url{https://github.com/SuperClaude-Org/SuperClaude_Framework}
\end{itemize}

%% --- Fin del contenido --- %%

\end{document}
